\documentclass[11pt,a4paper]{article}
\usepackage[margin=2.2cm]{geometry}
\usepackage[T1]{fontenc}
\usepackage[utf8]{inputenc}
\usepackage[provide=*,french]{babel}
\usepackage{lmodern,microtype}
\usepackage{xcolor}
\usepackage{enumitem}
\usepackage{fancyhdr}
\usepackage[hidelinks]{hyperref}
\definecolor{accent}{HTML}{245B65}
\hypersetup{pdftitle={Community Dictionaries : construire ensemble des ressources pour sa langue},pdfauthor={ChatGPT C-LARA-Instance et Manny Rayner},pdfsubject={Note pour les partenaires des communautés autochtones, première version, 5 octobre 2026}}
\setlength{\parindent}{0pt}
\setlength{\parskip}{6pt}
\setlength{\emergencystretch}{2em}
\setlist[itemize]{leftmargin=1.4em,itemsep=3pt,topsep=3pt}
\pagestyle{fancy}
\fancyhf{}
\fancyhead[L]{\small\textcolor{accent}{Community Dictionaries}}
\fancyhead[R]{\small Point d'étape --- octobre 2026}
\fancyfoot[C]{\small\thepage}
\setlength{\headheight}{14pt}
\renewcommand{\headrulewidth}{0.3pt}
\newcommand{\rubrique}[1]{\vspace{5pt}{\large\bfseries\textcolor{accent}{#1}\par}\vspace{2pt}}
\begin{document}
\thispagestyle{plain}
\begin{center}
{\LARGE\bfseries\textcolor{accent}{Community Dictionaries}\par}
\vspace{5pt}
{\Large Construire ensemble des ressources pour sa langue\par}
\vspace{8pt}
ChatGPT C-LARA-Instance et Manny Rayner --- Projet C-LARA\\[4pt]
{\small 5 octobre 2026 · Première version pour discussion}\\[7pt]
{\small\itshape Le texte intégral de cette note a été rédigé par ChatGPT
C-LARA-Instance,\\ à partir de ses échanges avec Manny Rayner.}
\end{center}

\textbf{Community Dictionaries est une application utilisable sur téléphone pour
construire, partager et utiliser des dictionnaires illustrés et sonores.} Son
point de départ est simple : les photographies, les voix et les connaissances des
personnes qui y participent. Une première version fonctionne déjà en ligne.
Nous souhaitons maintenant la confronter aux pratiques et aux attentes des
communautés, avec leurs partenaires linguistes et pédagogues.

\rubrique{Partir de la vie quotidienne}
On peut photographier un objet, un animal ou une scène, puis lui associer un mot
ou une expression, une traduction et un enregistrement. Il est également possible
de commencer par la voix ou par le texte. Une personne apporte une photographie,
une autre enregistre la langue : nul besoin que chacun fasse tout. Des demandes
entre partenaires et des commentaires permettent de compléter et de discuter les
entrées, avant leur validation par les personnes chargées de la révision.

Le premier dictionnaire développé par Cathy et Manny compte 61 entrées en
suédois, avec des traductions anglaises. Cathy a choisi des photographies prises
à la maison, en promenade ou lors d'occasions antérieures ; Manny a ajouté les
mots et les enregistrements suédois. Le contenu renvoie ainsi à des expériences
vécues, et pas seulement à une liste de vocabulaire préparée à l'avance.

Une même image peut évoquer plusieurs choses. Une photographie d'un chat sur un
canapé peut servir pour \emph{chat} comme pour \emph{canapé}. L'application permet
de relier manuellement plusieurs mots à une image et plusieurs images à un mot.
On peut parcourir le dictionnaire par images ou par mots, écouter un mot et
révéler sa traduction sans quitter l'image. Cette souplesse permet de discuter
ce que l'on veut montrer, plutôt que d'imposer une seule lecture de la scène.

\rubrique{Un outil façonné par les échanges avec les communautés}
Les échanges avec Sophie Rendina, nourris par son travail à Kowanyama, ont orienté
plusieurs choix centraux : un usage simple sur téléphone, la collaboration autour
d'images et de voix, des jeux avec indices visuels, la possibilité d'expérimenter
avec des images générées, et surtout le contrôle de ses contributions. Les essais
de Cathy et Manny ont également conduit à simplifier les commandes et les parcours.

Les communautés pourront conserver ce qui leur convient et demander des
ajustements précis. L'objectif est de rendre quelques activités utiles faciles à
réaliser, tout en évitant que l'ajout de possibilités rende l'application difficile
à apprendre. Le confort des personnes peu à l'aise avec l'écrit reste un critère
à vérifier avec elles.

\newpage
\rubrique{Utiliser les images et les voix pour apprendre}
Les ressources du dictionnaire alimentent déjà plusieurs activités :
\begin{itemize}
\item \textbf{Des cartes de révision}, avec réponses à choix multiple ou réponse
à révéler. Elles associent image, texte et audio dans les six directions possibles :
par exemple, entendre un mot et choisir son image, ou voir une image et choisir
l'enregistrement correspondant.
\item \textbf{Des mots croisés avec indices illustrés}, construits à partir des
mots et des images du dictionnaire.
\item \textbf{Des mots mêlés avec indices illustrés}, où l'on retrouve dans une
grille les mots évoqués par les images.
\end{itemize}
Les activités utilisent le contenu validé et peuvent être limitées à une catégorie.
Elles ne nécessitent pas de génération par IA pendant le jeu. La possibilité de
travailler avec ses propres images et les voix de son entourage nous paraît
prometteuse pour l'apprentissage ; son effet sur la motivation et les acquis reste
à évaluer en situation réelle.

\rubrique{Des aides par IA, lorsqu'elles sont souhaitées}
La photographie et l'enregistrement humain constituent un parcours complet, même
pour une langue que les modèles connaissent mal ou pas du tout. Des outils
supplémentaires sont disponibles lorsque le groupe souhaite les utiliser :
\begin{itemize}
\item \textbf{Créer des illustrations.} Une personne habilitée définit d'abord un
style commun, examine une image d'essai et approuve le style. Les illustrations
suivantes reprennent cette description ; chacune peut être examinée avant d'être
ajoutée. Cette fonction est désactivée par défaut et ne remplace pas l'appareil photo.
\item \textbf{Apprendre à partir d'une photo.} Pour les langues prises en charge,
l'IA propose l'identification d'un objet principal, un mot ou une expression dans
la langue étudiée et, si souhaité, une voix de synthèse. La proposition est
confirmée ou corrigée avant d'entrer dans le dictionnaire.
\item \textbf{Créer une version dans une autre langue.} Les images sont conservées ;
les textes sont adaptés en tenant compte du texte d'origine et de l'image, puis
un nouvel audio peut être produit. Le dictionnaire d'origine est préservé. Une
estimation du coût précède l'accord de lancement ; les résultats sont révisables.
\end{itemize}
Ces aides impliquent des services externes et doivent respecter les choix du groupe
sur les données qui peuvent leur être transmises. Elles peuvent faire des erreurs,
notamment de prononciation ou d'interprétation culturelle : le jugement des
locuteurs et des personnes concernées reste déterminant.

\rubrique{Garder la maîtrise de ses contributions}
Chaque dictionnaire est privé et accessible sur invitation. Un membre peut, après
confirmation, retirer toutes ses contributions. Elles lui restent accessibles en
privé ; jusqu'à leur restauration, il ne peut plus parcourir le dictionnaire
partagé ni y contribuer. Une commande lui permet de revenir et de restaurer son
contenu, en conservant les attributions. Si la personne responsable se retire,
elle organise une relève ; sans successeur, le dictionnaire est mis en archive.

Le service est actuellement hébergé par le projet sur AWS, avec un accès technique
d'administration. Le retrait bloque le partage dans l'application ; il n'efface
ni les sauvegardes ni les copies déjà téléchargées ou transmises à un service
externe. Hébergement, conservation, réutilisation et restrictions culturelles
restent à convenir avec chaque communauté. Ces commandes concrétisent une partie
du contrôle ; la souveraineté des données demande aussi des accords clairs.

\newpage
\rubrique{Un développement rapide, avec des retours humains essentiels}
Une particularité du projet est sa méthode de développement. L'assistant d'IA
utilisé pour ces travaux, GPT-6 Astra en mode Extra High, a rédigé le code, les
tests et la documentation. Manny a installé les versions successives, effectué
les essais et transmis ses observations et celles des autres participants.
L'application réutilise également l'infrastructure et plusieurs fonctions de
C-LARA-2, notamment pour les médias et les jeux.

Selon l'estimation rétrospective de Manny, le travail de conception et de
programmation de l'assistant représente \textbf{environ cinq heures, réparties
sur environ une semaine d'échanges et d'itérations}. Ce chiffre n'inclut pas
l'ensemble du temps humain consacré aux installations, aux essais et aux échanges,
ni le développement antérieur de C-LARA-2. Il donne un ordre de grandeur, sans
constituer une mesure chronométrée ou une comparaison contrôlée entre modèles.
Manny constate une nette accélération par rapport à ses expériences antérieures.

L'intérêt pratique est de pouvoir partir d'une demande exprimée simplement,
l'essayer rapidement, puis l'ajuster. Par exemple, les échanges sur le retrait des
contributions ont conduit à remplacer plusieurs commandes détaillées par une
logique plus facile à comprendre : \emph{se retirer en conservant son contenu,
puis pouvoir revenir}. Les essais ont aussi permis de rapprocher les commandes
de sauvegarde des actions auxquelles elles se rapportent et d'accélérer la
révision des versions linguistiques.

\rubrique{Ce qui a été essayé, et ce qu'il reste à apprendre}
La création et l'utilisation d'entrées ont été essayées sur ordinateur et sur un
téléphone réel. Les versions récentes sont installées sur le serveur du projet.
Le retrait et la restauration des contributions ont été testés, ainsi que la
génération d'images et les activités de révision. Des essais automatisés complètent
ces vérifications. Axel a aussi rapporté une première utilisation positive pour
l'islandais : photographier des objets, confirmer leur identification et écouter
leur nom dans cette langue.

Le dictionnaire suédois de 61 entrées a maintenant été adapté en français puis en
allemand. Manny juge les deux résultats très positifs. Il a toutefois repéré
plusieurs prononciations françaises insatisfaisantes et un problème avec
\emph{curry} en allemand ; sa connaissance moins assurée de l'allemand peut en
avoir laissé passer d'autres. Ces essais montrent un fonctionnement utile, sans
établir un taux de justesse ni valider ces aides pour les langues autochtones.

\rubrique{Prochaine étape : essayer et choisir ensemble}
Sophie prévoit une présentation à Kowanyama au début de novembre. Cette note peut
aussi servir de base aux échanges avec les collègues de Nouméa, si une
expérimentation autour des langues kanak les intéresse. Les premiers retours
pourraient porter sur des questions très concrètes : quelles activités donnent
envie de revenir ? Quels mots et boutons sont peu clairs ? Qui souhaite contribuer,
réviser ou utiliser les aides par IA ? Quelles règles de partage faut-il respecter ?

À plus long terme, nous souhaitons réduire progressivement le besoin d'un
intermédiaire technique pour adapter et maintenir l'outil. La rapidité actuelle
encourage cette recherche ; une maintenance entièrement autonome n'est pas encore
démontrée. Les communautés et leurs partenaires doivent continuer à définir les
besoins, les règles et les critères d'un résultat utile.

\vfill
{\small Application : \url{https://c-lara-2.c-lara.org/community-dictionaries/}\\
Code et documentation ouverts : \url{https://github.com/mannyrayner/C-LARA-2}\\
Cette note décrit la situation au 5 octobre 2026 et invite à la discussion.}
\end{document}
